\section{Conclusion}\label{sec:concl}
We investigated multi-agent coordination and the value of agent specialization within an orchestrated technical interview preparation system. By introducing specialized Strategy and Coaching agents with private state ledgers, typed proposal contracts, and unidirectional peer messaging, we demonstrated that stateful coordination alters difficulty trajectories in all 160 replay sessions, smoothing level changes ($-2.67$) and reversals ($-1.88$) relative to a single-policy PPO baseline, while adding $+0.59$ proactive probe turns driven by specific conceptual weakness clusters.

Crucially, our empirical and architectural analysis challenges the assumption that agent specialization inherently optimizes decision quality. Under deterministic heuristic rules, an ablated coaching condition and an equivalent unified agent produce mathematically identical difficulty trajectories ($\mathbb{P}(\text{agree}) = 1.0$). The genuine value of agent specialization lies not in algorithmic optimization, but in software engineering properties: disjoint state ledgers, independent auditability, and bulkhead fault containment, balanced against a $+1.04$\,ms median dispatch latency overhead. By isolating coordination mechanisms via threshold decomposition and cluster-bootstrap replay, this work provides a reproducible, evidence-based foundation for designing bounded-authority multi-agent systems in adaptive educational computing.
